\section{Physical AI: la combinación de sistemas de Omniverse y Cosmos}

\subsection{Del sentido común lingüístico al sentido común físico}

Jensen expone la dificultad de Physical AI de manera muy directa: los robots no solo deben «saber hablar», sino también «entender cómo se mueve el mundo». Los modelos de lenguaje resuelven la inferencia semántica, mientras que los sistemas robóticos necesitan dominar además las restricciones físicas, las relaciones espaciales y la retroalimentación de ejecución.

\begin{importantbox}{Physical AI en una frase}
\textit{``Everything that moves will be robotic someday.''}

En términos de ingeniería, esta frase se traduce en tres capacidades: percibir (Perception), decidir (Planning) y ejecutar (Control). Si cualquiera de estos eslabones es inestable, el sistema no puede llegar a implementarse a escala.
\end{importantbox}

\subsection{Por qué construir primero un «campo de entrenamiento virtual escalable»}

En la entrevista se menciona que entrenar en el mundo real es costoso y lento, y además el hardware robótico se desgasta. El valor de Omniverse está en ofrecer un entorno físico 3D controlable, reproducible y paralelizable; el valor de Cosmos está en dotar a esos entornos de conocimientos previos del mundo más ricos y de capacidad de generación de escenas.

\begin{knowledgebox}{La división de funciones entre Omniverse y Cosmos}
\begin{itemize}
    \item Omniverse: motor de simulación y banco de trabajo de gemelos digitales, encargado de la «consistencia física» y la «ejecución del entrenamiento».
    \item Cosmos: capacidad de modelo del mundo, encargada de la «diversidad de escenas» y el «puente entre semántica y física».
\end{itemize}
El objetivo de combinar ambos es transformar el entrenamiento de robots de un «ensayo y error costoso» a una «iteración industrializable».
\end{knowledgebox}

\begin{warningbox}{Error común: creer que basta con tener un world model para sustituir la simulación directamente}
El world model y la simulación no están en una relación de sustitución. El primero es bueno para generar y generalizar; la segunda, para restringir y validar. Los sistemas robóticos necesitan que coexistan tres elementos: «capacidad de generación + consistencia física + evaluación de circuito cerrado»; ninguno puede faltar.
\end{warningbox}

\subsection{Resumen de la sección}
Lo esencial de Physical AI no es el desempeño de un modelo puntual, sino la infraestructura de entrenamiento. Omniverse + Cosmos representan la «fábrica de datos + fábrica de conocimientos previos de la era de los robots».

